\chapter{系统概述}\label{chap:overview}

第三章已把历史学家在档案重建中的认知负担抽象为 DG1--DG6 六项设计目标。本章不展开算法细节，仅从高层视角说明本系统如何回应这些目标：后端承担哪些计算工作、前端承担哪些审议任务，二者共同遵循的可视分析原则。图建模与协商协议见第五章，视图与交互细节见第六章。

\section{系统目标与用户需求映射}\label{sec:overview-needs}

DG1--DG6 看似分散，实则可归并为三类高层需求。其一为\textbf{可解释性}：DG1 要求把图模型捕捉的拓扑模式抽取成可命名的结构证据，DG4 要求专家在协商进行时可注入领域知识，二者共同避免标量分数遮蔽史学语境。其二为\textbf{可审计性}：DG2 要求每条被接受的边附带来源标记并保留多轮协商日志，使方法论审稿人可独立复核\citep{amershi2014modeltracker}。其三为\textbf{可逆性与跨队列规则迁移}：DG5 要求一键回退后界面与后台状态一致；DG3 与 DG6 则要求在昂贵推理前提供过滤入口，并把整定的偏好沉淀为可在新队列复用的空间先验。

回应上述三类需求的难点在于：任何单一证据线都不足以兼顾解释、审计与迁移。仅靠图嵌入的标量分数无法解释拓扑来源；仅靠多智能体协商虽给出文本论证，但若缺乏图结构锚定便有事后合理化之嫌\citep{turpin2023faithfulness}；纯人工逐对审议则难以扩展到大规模队列。本系统因此采取"候选排序、结构指纹、协商论证与人在环可视分析"四条证据流协同的路线，并以混合主动\citep{green2009mixedinitiative} 与知识生成型可视分析理论\citep{sacha2014knowledge} 为指导原则。系统总体架构如图~\ref{fig:framework}所示，分校准与部署两阶段，把单次队列的方法学投入沉淀为跨队列可迁移的先验。

\begin{figure}[!htbp]
    \centering
    \includegraphics[width=0.95\linewidth]{Framework}
    \caption{GeneaLink 两阶段框架。校准阶段在带有真值标注的队列上以异构图变换器候选排序与子图模式提取为基础，由历史学家在协商竞技场内整定权重并落盘为单元格层级先验；部署阶段对新队列重做候选生成，并自动加载先验注入协商首轮提示。}
    \label{fig:framework}
\end{figure}

\section{后端引擎与前端界面分工}\label{sec:overview-arch}

后端引擎承担四件相互衔接的工作。其一是\textbf{异构图建模}：把历史档案中的人物、家户、社区、八旗及其多类亲缘与共居关系组织为带时间戳的异构图，并以时间因果约束防止未来信息泄漏。其二是\textbf{候选评分}：以异构图变换器为编码器，对每对候选给出排序分数，使下游推理能先聚焦在高置信子集。其三是\textbf{结构指纹提取}：对候选周边的局部拓扑做模式抽取，把直系同胞、共父、母方交表等社会学常识结构以可命名的形式呈现，承担 DG1 的结构证据职责。其四是\textbf{多智能体协商}：代表丈夫与妻子两侧的智能体在六轮结构化对话中引用上述结构指纹与生活事件证据，给出最终评分与可逐轮回放的论证日志，呼应生成式智能体的协商范式\citep{park2023generative}。后端的产出包括候选排序、结构指纹、协商分数与日志四类，每类都附带来源标记。

前端界面把上述产出以六个相互联动的视图呈现给历史学家：性能概览跟踪队列层级的进展曲线，流程视图以账本记录每条已接受边的来源与时间，嵌入视图以稳定的二维空间承载候选并支持空间过滤，二分图视图支持单对审视与批量接受，协商竞技场可视化分轮辩论并接受提示注入，规则视图允许整定权重并把偏好落盘为跨队列可复用的先验。所有视图共享同一选区与来源色谱，操作以一次动画帧延迟传播到其余视图。

由此形成清晰的分工：后端把图结构与协商证据编织为带审计线索的多层证据流，前端提供提示注入与一键回滚两条反向通道，使专家既能在协商中干预论证，又能事后撤销任何已接受的边，从而共同保证系统状态可审计、可逆且可迁移。

\section{可视化设计规则}\label{sec:overview-vis-rules}

在前后端协同的框架下，本系统恪守四条可视分析原则，使 DG1--DG6 在每处交互中都不被牺牲。

\textbf{语义交互优先于装饰交互}。任一鼠标事件都映射为可叙事化的史学动作，如"接受这一对"或"对该智能体降权"，而非孤立的图形操作\citep{endert2017state}。该规则把 DG1 的结构证据延伸到交互层，使专家在视图间切换不丢失语义上下文。

\textbf{来源可见且不可篡改}。所有进入家谱的边都显式携带来源色条，标记其来自模型、协商或专家录入；多轮协商日志一经写入即只读保存。该规则直接服务 DG2 与 DG3，使审计跟踪在视觉层面具备稳定锚点。

\textbf{反向操作不丢失状态}。任何提交都配套一键回滚通道，回滚后候选回到候选池并保持嵌入坐标稳定，关联视图同步复位。该规则旨在消除"界面假性成功而后台仍持有旧状态"的隐性裂缝，是 DG5 的工程落点。

\textbf{跨视图一次动画帧响应}。任一视图发起的选区变化都在下一帧内更新到其余视图。该规则与前述三条共同构成混合主动协作的工作面，使专家既能以批量节奏推进高置信子集，又能深入审议边界对，并让校准过的偏好经空间先验自动迁移到新队列，呼应 DG6 的跨队列复用诉求。
